% GENERATED FILE. Edit canonical lessons, then run npm run generate:books.
% canonical-source-hash: fnv1a64:89d2827d7f653edb
% canonical-lessons: PT-C227-papelaria, PT-C227-livraria, PT-C227-peixaria, PT-C227-centro-comercial, PT-C227-montra

\chapter{Stationer's, Bookshop, Fishmonger's, Shopping centre, Shop window}
\label{ch:pt-stationer-s-bookshop-fishmonger-s-shopping-centre-shop-window}
\hlchaptermodality{227}

\begin{chapteropening}
\textbf{By the end of this chapter:} \emph{I can say a stationer's, a bookshop, a fishmonger's, a shopping centre and a shop window.}

Complete the last lesson of chapter 227: I can say a stationer's, a bookshop, a fishmonger's, a shopping centre and a shop window.
\end{chapteropening}

\section[a papelaria]{a papelaria — a stationer's}

\label{lesson:PT-C227-papelaria}



\begin{quote}
Before the new one: say the Portuguese for tennis, then the Portuguese for a victory.
\end{quote}

\subsection*{You'll want to know: a papelaria}
\textbf{a papelaria} — \textquotedblleft{}a stationer's\textquotedblright{}.

\begin{cousinweb}
\textbf{papelaria} is built from \emph{papel}, \textquotedblleft{}paper\textquotedblright{}, and \emph{-aria}, \textquotedblleft{}shop\textquotedblright{}.
\end{cousinweb}

\begin{grammarlens}[title={What you've built}]
One more word for this chapter.
\end{grammarlens}

\subsection*{Your turn}
\begin{itemize}
\raggedright
  \item \emph{Say it:} \emph{a papelaria}
  \item \emph{Say it:} \emph{a papelaria}, once more
  \item \emph{Say it:} say \emph{a papelaria}
  \item \emph{Recall:} say the Portuguese for tennis, then the Portuguese for a victory, then say \emph{a papelaria} again
\end{itemize}

\begin{tcolorbox}[breakable,colback=teal!4,colframe=teal!35!black,title={Before you move on}]
What does a papelaria mean? (\textquotedblleft{}A stationer's\textquotedblright{}.) Say it once more.
\end{tcolorbox}
\section[a livraria]{a livraria — a bookshop}

\label{lesson:PT-C227-livraria}



\begin{quote}
Before the new one: say the Portuguese for a victory, then the Portuguese for a stationer's.
\end{quote}

\subsection*{You'll want to know: a livraria}
\textbf{a livraria} — \textquotedblleft{}a bookshop\textquotedblright{}.

\begin{cousinweb}
\textbf{livraria} is built from \emph{livro}, \textquotedblleft{}book\textquotedblright{} ($\leftarrow$ Latin \emph{liber}, \textquotedblleft{}book\textquotedblright{}), and \emph{-aria}, \textquotedblleft{}shop\textquotedblright{}. It is not a library.
\end{cousinweb}

\begin{grammarlens}[title={What you've built}]
One more word for this chapter.
\end{grammarlens}

\subsection*{Your turn}
\begin{itemize}
\raggedright
  \item \emph{Say it:} \emph{a livraria}
  \item \emph{Say it:} \emph{a livraria}, once more
  \item \emph{Say it:} say \emph{a livraria}
  \item \emph{Recall:} say the Portuguese for a victory, then the Portuguese for a stationer's, then say \emph{a livraria} again
\end{itemize}

\begin{tcolorbox}[breakable,colback=teal!4,colframe=teal!35!black,title={Before you move on}]
What does a livraria mean? (\textquotedblleft{}A bookshop\textquotedblright{}.) Say it once more.
\end{tcolorbox}
\section[a peixaria]{a peixaria — a fishmonger's}

\label{lesson:PT-C227-peixaria}



\begin{quote}
Before the new one: say the Portuguese for a stationer's, then the Portuguese for a bookshop.
\end{quote}

\subsection*{You'll want to know: a peixaria}
\textbf{a peixaria} — \textquotedblleft{}a fishmonger's\textquotedblright{}.

\begin{cousinweb}
\textbf{peixaria} is built from \emph{peixe}, \textquotedblleft{}fish\textquotedblright{}, and \emph{-aria}, \textquotedblleft{}shop\textquotedblright{}.
\end{cousinweb}

\begin{grammarlens}[title={What you've built}]
One more word for this chapter.
\end{grammarlens}

\subsection*{Your turn}
\begin{itemize}
\raggedright
  \item \emph{Say it:} \emph{a peixaria}
  \item \emph{Say it:} \emph{a peixaria}, once more
  \item \emph{Say it:} say \emph{a peixaria}
  \item \emph{Recall:} say the Portuguese for a stationer's, then the Portuguese for a bookshop, then say \emph{a peixaria} again
\end{itemize}

\begin{tcolorbox}[breakable,colback=teal!4,colframe=teal!35!black,title={Before you move on}]
What does a peixaria mean? (\textquotedblleft{}A fishmonger's\textquotedblright{}.) Say it once more.
\end{tcolorbox}
\section[o centro comercial]{o centro comercial — a shopping centre}

\label{lesson:PT-C227-centro-comercial}



\begin{quote}
Before the new one: say the Portuguese for a bookshop, then the Portuguese for a fishmonger's.
\end{quote}

\subsection*{You'll want to know: o centro comercial}
\textbf{o centro comercial} — \textquotedblleft{}a shopping centre\textquotedblright{}.

\begin{cousinweb}
\textbf{centro comercial} is \textquotedblleft{}commercial centre\textquotedblright{}: \emph{centro} $\leftarrow$ Latin \emph{centrum} (EtymDB 2.1), and \emph{comercial} $\leftarrow$ Late Latin \emph{commercialis}, from \emph{commercium}, \textquotedblleft{}trade\textquotedblright{} (EtymDB 2.1).
\end{cousinweb}

\begin{grammarlens}[title={What you've built}]
One more word for this chapter.
\end{grammarlens}

\subsection*{Your turn}
\begin{itemize}
\raggedright
  \item \emph{Say it:} \emph{o centro comercial}
  \item \emph{Say it:} \emph{o centro comercial}, once more
  \item \emph{Say it:} say \emph{o centro comercial}
  \item \emph{Recall:} say the Portuguese for a bookshop, then the Portuguese for a fishmonger's, then say \emph{o centro comercial} again
\end{itemize}

\begin{tcolorbox}[breakable,colback=teal!4,colframe=teal!35!black,title={Before you move on}]
What does o centro comercial mean? (\textquotedblleft{}A shopping centre\textquotedblright{}.) Say it once more.
\end{tcolorbox}
\section[a montra]{a montra — a shop window}

\label{lesson:PT-C227-montra}



\begin{quote}
Before the new one: say the Portuguese for a fishmonger's, then the Portuguese for a shopping centre.
\end{quote}

\subsection*{You'll want to know: a montra}
\textbf{a montra} — \textquotedblleft{}a shop window\textquotedblright{}.

\begin{cousinweb}
\textbf{montra} $\leftarrow$ French \emph{montre}, \textquotedblleft{}display\textquotedblright{} (EtymDB 2.1) — borrowed. The French word comes from \emph{montrer}, \textquotedblleft{}to show\textquotedblright{}.
\end{cousinweb}

\begin{grammarlens}[title={What you've built}]
One more word for this chapter.
\end{grammarlens}

\subsection*{Your turn}
\begin{itemize}
\raggedright
  \item \emph{Say it:} \emph{a montra}
  \item \emph{Say it:} \emph{a montra}, once more
  \item \emph{Say it:} say \emph{a montra}
  \item \emph{Recall:} say the Portuguese for a fishmonger's, then the Portuguese for a shopping centre, then say \emph{a montra} again
\end{itemize}

\begin{tcolorbox}[breakable,colback=teal!4,colframe=teal!35!black,title={Before you move on}]
What does a montra mean? (\textquotedblleft{}A shop window\textquotedblright{}.) Say it once more.
\end{tcolorbox}
